\documentclass[11pt]{article}
\newcommand{\DocShort}{Cronograma}
\usepackage{fontspec}
\usepackage[spanish,es-noshorthands]{babel}
\decimalpoint
\usepackage{geometry}
\geometry{letterpaper, top=2.5cm, bottom=2.5cm, left=2.8cm, right=2.8cm}
\usepackage[expansion=false]{microtype}
\usepackage{parskip}
\setlength{\parskip}{6pt}
\usepackage{amsmath,amssymb}
\usepackage{booktabs}
\usepackage{array}
\usepackage{longtable}
\usepackage{caption}
\usepackage[dvipsnames,table]{xcolor}
\usepackage{enumitem}
\setlist[itemize]{topsep=2pt, itemsep=2pt, parsep=0pt}
\setlist[enumerate]{topsep=2pt, itemsep=2pt, parsep=0pt}
\usepackage{fancyhdr}
\usepackage[hidelinks,pdfencoding=auto]{hyperref}
\usepackage{titlesec}
\usepackage{tcolorbox}
\tcbuselibrary{breakable}
\usepackage{pdfpages}

\definecolor{darkblue}{HTML}{1B3A5C}
\definecolor{medblue}{HTML}{2E6DA4}
\definecolor{lightblue}{HTML}{D6E8F7}
\definecolor{teal}{HTML}{1A7A6E}
\definecolor{lightteal}{HTML}{D0EDEA}
\definecolor{lightgrey}{HTML}{F5F5F5}
\definecolor{midgrey}{HTML}{6F6F6F}
\definecolor{warnyellow}{HTML}{FFF3CD}

\titleformat{\section}
{\large\bfseries\color{darkblue}}
{\thesection}{0.75em}{}[{\color{medblue}\titlerule[0.8pt]}]
\titleformat{\subsection}
{\normalsize\bfseries\color{darkblue}}
{\thesubsection}{0.75em}{}
\titleformat{\subsubsection}
{\normalsize\bfseries\color{medblue}}
{\thesubsubsection}{0.75em}{}

\pagestyle{fancy}
\fancyhf{}
\fancyhead[L]{\small\color{midgrey}ICV 442 Acueductos y Alcantarillados}
\fancyhead[R]{\small\color{midgrey}\DocShort}
\fancyfoot[C]{\small\color{midgrey}\thepage}
\renewcommand{\headrulewidth}{0.4pt}
\renewcommand{\footrulewidth}{0pt}

\rowcolors{2}{lightgrey}{white}
\renewcommand{\arraystretch}{1.18}

\newtcolorbox{infobox}{
  breakable,
  colback=lightblue,
  colframe=medblue,
  boxrule=0.8pt,
  arc=3pt,
  left=8pt,
  right=8pt,
  top=7pt,
  bottom=7pt
}

\newtcolorbox{warnbox}{
  breakable,
  colback=warnyellow,
  colframe=orange!70!black,
  boxrule=0.8pt,
  arc=3pt,
  left=8pt,
  right=8pt,
  top=7pt,
  bottom=7pt
}

\newcommand{\doccover}[3]{
  \begin{center}
  {\small\color{midgrey}Pontificia Universidad Cat\'olica Madre y Maestra\\Escuela de Ingenier\'ia Civil y Ambiental}

  \vspace{1.0cm}
  {\color{medblue}\rule{\linewidth}{2pt}}
  \vspace{0.3cm}

  {\Large\bfseries\color{darkblue}ICV-442-T Acueductos y Alcantarillado}\\[0.35em]
  {\LARGE\bfseries\color{darkblue}#1}\\[0.25em]
  {\large\itshape #2}

  \vspace{0.3cm}
  {\color{medblue}\rule{\linewidth}{2pt}}
  \vspace{0.9cm}
  \end{center}

  \begin{center}
  \begin{tabular}{ll}
  \toprule
  \textbf{Asignatura} & ICV-442-T Acueductos y Alcantarillado \\
  \textbf{Profesor} & Ricardo Rafael Hern\'andez Moreira, PhD \\
  \textbf{Periodo acad\'emico} & Mayo--agosto de 2026 \\
  \textbf{Proyecto} & Dise\~no de sistemas de saneamiento para Las Charcas, Azua \\
  \textbf{Versi\'on} & #3 \\
  \bottomrule
  \end{tabular}
  \end{center}

  \vspace{0.4cm}
}


\begin{document}
\doccover{Cronograma del proyecto}{Secuencia de hitos y entregas}{18 de mayo de 2026}

\begin{infobox}
La din\'amica general del curso combina avances breves en clase los jueves y entregas por plataforma los viernes. El cronograma est\'a pensado para obligar a trabajar de forma incremental y evitar que la carga se desplace a las \'ultimas semanas.
\end{infobox}

\section{Calendario de hitos}
\small
\begin{longtable}{p{0.08\linewidth} p{0.17\linewidth} p{0.33\linewidth} p{0.22\linewidth} p{0.08\linewidth}}
\toprule
\textbf{Sem.} & \textbf{Fechas} & \textbf{Trabajo principal} & \textbf{Entrega del viernes} & \textbf{Peso} \\
\midrule
\endfirsthead
\toprule
\textbf{Sem.} & \textbf{Fechas} & \textbf{Trabajo principal} & \textbf{Entrega del viernes} & \textbf{Peso} \\
\midrule
\endhead
3 & 18--22 mayo & Presentaci\'on del proyecto, conformaci\'on de equipos, revisi\'on de datos disponibles y organizaci\'on interna. & Acta de conformaci\'on del equipo y plan interno de trabajo. & 0\% \\
4 & 25--29 mayo & Proyecci\'on poblacional, delimitaci\'on preliminar del \'area de servicio y revisi\'on de fuentes de datos. & Memoria de proyecci\'on poblacional. & 3\% \\
5 & 1--5 junio & Definici\'on de par\'ametros base, criterios normativos y brechas de informaci\'on. & Tabla maestra de par\'ametros y brechas. & 2\% \\
6 & 8--12 junio & Trazado preliminar y compatibilidad entre las tres redes. & Anteproyecto con esquemas de trazado. & 5\% \\
7 & 15--19 junio & Configuraci\'on y primera corrida del modelo de agua potable. & Reporte preliminar EPANET. & 4\% \\
8 & 22--26 junio & Configuraci\'on y primera corrida de los modelos sanitario y pluvial. & Reporte preliminar SWMM. & 4\% \\
9 & 29 junio--3 julio & Integraci\'on de resultados y redacci\'on de memoria. & Borrador integrado del proyecto. & 2\% \\
10 & 6--10 julio & Revisi\'on por pares y correcciones dirigidas. & Informe de revisi\'on por pares. & 0\% \\
11 & 13--17 julio & Ajuste fino del documento y ensayo de presentaci\'on. & Versi\'on preliminar para revisi\'on docente. & 0\% \\
12 & 20--24 julio & Cierre documental, verificaci\'on de formato y preparaci\'on de anexos. & Entrega final del proyecto escrito antes del viernes 24 de julio de 2026, 23:59. & 20\% \\
13 & 27--31 julio & Preparaci\'on de la exposici\'on y ajustes de material visual. & Sin entrega formal. & --- \\
14 & 3--4 agosto & Presentaciones orales finales. & Exposici\'on y defensa del proyecto. & 5\% \\
\bottomrule
\end{longtable}
\normalsize

\section{R\'egimen de consecuencias}
\begin{itemize}
\item Los descuentos por retraso se aplican sobre la ponderaci\'on del hito correspondiente.
\item Una entrega con m\'as de tres d\'ias de retraso no se califica, salvo autorizaci\'on expresa y documentada.
\item La revisi\'on por pares de la semana 10 es requisito habilitante para aceptar la entrega final.
\item La entrega final no suma un 20\% adicional a los hitos previos; consolida y valida el componente escrito del proyecto.
\end{itemize}

\section{Uso del cronograma}
\begin{itemize}
\item Los equipos deben convertir este calendario en un cronograma interno con responsables y fechas m\'as cortas.
\item Las presentaciones breves de los jueves deben reportar avances reales, dificultades y decisiones para la semana siguiente.
\item El cronograma no est\'a dise\~nado para premiar el trabajo acumulado al final, sino para sostener progreso continuo.
\end{itemize}

\end{document}
